\section{Автоморфизмы расширений}

\begin{definition}
    Пусть $F \subset E$ --- расширение полей. \emph{$F$-автоморфизм} поля $E$ --- это автоморфизм $\varphi\colon E \to E
    $, тождественный на $F$: $\varphi(a) = a$ для всех $a \in F$.

    Множество всех $F$-автоморфизмов $E$ обозначается $\mathrm{Aut}(E/F)$ и является группой относительно композиции --- \emph{группой Галуа}.
\end{definition}


\begin{examplelist}
    \item $\mathrm{Aut}(\mathbb{C}/\mathbb{R})$ состоит из двух элементов: $\mathrm{id}$ и комплексное сопряжение
    $\overline{\phantom{z}}$.

    \item $E = \mathbb{C}(x)$ над $\mathbb{C}$. Любой $\mathbb{C}$-автоморфизм $\varphi$ определяется значением
    $\varphi(x) = u \in E$. По следствию из леммы Люрота $u = \frac{ax + b}{cx + d}$, $\begin{vmatrix}
    a & b \\ c & d \end{vmatrix} \neq 0$. Эта группа называется $\mathrm{PGL}_2(\mathbb{C})$ (проективная линейная группа).

    \item $E = \mathbb{C}(x_1, \dots, x_n)$ над $\mathbb{C}$ --- бесконечная группа автоморфизмов.
\end{examplelist}


\section{Порядок группы Галуа}

\begin{stmt}
    Пусть $f \in F[x]$ --- сепарабельный многочлен, $E$ --- поле разложения $f$. Тогда
    \[
        |\mathrm{Aut}(E/F)| = [E : F].
    \]

    \begin{proof}
        Так как $f$ сепарабелен, у него ровно $\deg f$ различных корней в $E$. Каждый $F$-гомоморфизм
        $\varphi\colon E \to E$ отправляет корни в корни. Но мы знаем, что $E$ --- конечное расширение
        $F$, поэтому $\varphi$ --- инъекция конечномерного пространства в себя, значит биекция, значит автоморфизм.
        Число таких гомоморфизмов $= [E : F]$ (так как $f$ сепарабелен).
    \end{proof}
\end{stmt}


\section{Неподвижное поле}

\begin{definition}
    Пусть $E$ --- поле, $G \leq \mathrm{Aut}(E)$ --- подгруппа автоморфизмов. \emph{Неподвижное поле}:
    \[
        E^G = \{e \in E \mid \sigma(e) = e\ \forall\, \sigma \in G\}.
    \]
\end{definition}

\begin{stmt}
    $E^G$ --- подполе $E$.
    \begin{proof}
        Очевидно: $0, 1 \in E^G$, и если $a, b \in E^G$, то $\sigma(a \pm b) = \sigma(a) \pm \sigma(b)
        = a \pm b$, $\sigma(ab) = \sigma(a)\sigma(b) = ab$, $\sigma(a^{-1}) = \sigma(a)^{-1} = a^{-1}$.
    \end{proof}
\end{stmt}


\section{Теорема Артина}

\begin{theorem}[Артина]
    Пусть $G$ --- конечная группа автоморфизмов поля $E$, $|G| = n$. Тогда
    \[
        [E : E^G] \leq |G|.
    \]

    \begin{proof}
        Положим $F \coloneqq E^G$, $G = \{g_1, g_2, \dots, g_n\}$. Хотим показать $\dim_F E \leq n$.

        Достаточно: для любых $m > n$ элементов $e_1, \dots, e_m \in E$ показать, что они линейно зависимы над $F$.

        Рассмотрим систему $n$ уравнений с $m$ неизвестными $x_1, \dots, x_m$:
        \[
            \begin{cases}
                g_1(e_1)\, x_1 + g_1(e_2)\, x_2 + \dots + g_1(e_m)\, x_m = 0 \\
                g_2(e_1)\, x_1 + g_2(e_2)\, x_2 + \dots + g_2(e_m)\, x_m = 0 \\
                \quad\vdots \\
                g_n(e_1)\, x_1 + g_n(e_2)\, x_2 + \dots + g_n(e_m)\, x_m = 0
            \end{cases}
        \]
        Так как $m > n$, существует ненулевое решение $(d_1, \dots, d_m) \in E^m$.

        Выберем решение с наименьшим числом ненулевых координат. Без ограничения общности $d_1 \neq 0
        $ и $d_i \in F$.

        Если какое-то $d_i \notin F$, то $\exists\, g_t \in G\colon g_t(d_i) \neq d_i$. Применим $g_t
        $ ко всем уравнениям системы: так как $\{g_t g_1, g_t g_2, \dots, g_t g_n\} = \{g_1, \dots, g_n\}
        $, новая система --- та же самая (с переставленными строками). Но решение новой системы --- $(g_t(d_1)
        , g_t(d_2), \dots, g_t(d_m))$.

        Два решения: $\Gamma_1 = (d_1, d_2, \dots, d_m)$ и $\Gamma_2 = (d_1, g_t(d_2), \dots, g_t(d_m)
        )$ (первая координата совпадает, так как $d_1 \in F$).

        Разность $\Gamma_1 - \Gamma_2$ --- тоже решение, у которого первая координата $= 0$ и строго меньше ненулевых координат.
        Противоречие с минимальностью.

        Значит все $d_i \in F$, и $e_1 d_1 + e_2 d_2 + \dots + e_m d_m = 0$ --- линейная зависимость над $F$.

        Итого $[E : F] \leq n = |G|$.
    \end{proof}
\end{theorem}
